\section{열역학, 온도미분과 열수지}\label{THM-CH2}

Chapter~\ref{EQ-CH1}의 평형해가 온도에 따라 이동하면 그 이동은
가역열과 연결된다. 반면 평형에서 벗어나 생긴 flux와 affinity의
곱은 비가역 entropy production이다. 두 항을 같은 경험적 온도보정으로
합치지 않는 것이 이 장의 핵심이다.

\subsection{고정상태, 고정전하와 경로미분}\label{THM-CH2-DERIVATIVES}

\paragraph{\physid{THM-001} 세 온도미분.}
\[
\left.\partial_T U_{\eq,j}\right|_{\xi_j},\qquad
\left.\partial_T V_{n,\mathrm{OCV}}\right|_q,\qquad
\left.\frac{\dd V_{\app}}{\dd T}\right|_{\rm path}
\]
를 분리한다. 마지막 항에는 상태변화, 분극, 저항과 속도론이 섞일
수 있으므로 reaction entropy coefficient가 아니다.

\paragraph{\physid{THM-002} 고정전하 OCV 미분.}
평형에서 식~\eqref{BAL-001-EQUATION}에
$\xi_j=\xi_{\eq,j}(V_n,T)$를 넣고 $q$를 고정해 미분한다.
$Q_\cell$, $a_j$와 기준량이 온도에 무관하다는 제한에서
\begin{equation}
\boxed{
\left.\frac{\partial V_{n,\mathrm{OCV}}}{\partial T}\right|_q
=-
\frac{
\left.\partial_T Q_\bg^\chem\right|_{V_n}
+\displaystyle\sum_j a_j
\left.\partial_T\xi_{\eq,j}\right|_{V_n}
}{
C_\bg^\chem
+\displaystyle\sum_j a_j
\left.\partial_{V_n}\xi_{\eq,j}\right|_T
}.
}
\label{THM-002-EQUATION}
\end{equation}
분자의 단위는 C K$^{-1}$, 분모의 단위는 C V$^{-1}$이므로
결과는 V K$^{-1}$이다. $a_j(T)$ 또는 기준용량의 온도의존을
허용하면 곱미분 항
$(\partial_Ta_j)(\xi_{\eq,j}-\xi_{j,0})$ 등을 분자에 추가한다.

\begin{validity}
식~\eqref{THM-002-EQUATION}의 분모가 0이 아니어야 implicit
function theorem을 적용할 수 있다. signed $a_j$를 쓰므로
$C_\bg^\chem>0$만으로 전체 분모 양수를 자동 보장하지 않는다.
선택한 평형 branch에서 total differential capacitance의 부호를
별도로 확인한다.
\end{validity}

\subsection{reaction entropy와 가역열}\label{THM-CH2-REVERSIBLE}

\paragraph{\physid{THM-003} 하나의 반응에서의 부호 유도.}
written forward reaction 하나에 대해
\begin{equation}
U_{\eq}=-\frac{\Delta_rG}{zF},\qquad
I_{\rm rxn}=zF\dot n_{\rm rxn}
\label{THM-003-CONVENTION}
\end{equation}
로 정의한다. $(\partial\Delta_rG/\partial T)_{\rm state}=-\Delta_rS$
이므로
\begin{equation}
\boxed{
\Delta_rS=zF
\left.\frac{\partial U_{\eq}}{\partial T}\right|_{\rm state}.
}
\label{THM-003-ENTROPY}
\end{equation}
control volume에 남는 열이 양인 규약에서는 가역 heat generation이
$-T\Delta_rS\dot n_{\rm rxn}$이므로
\begin{equation}
\boxed{
\dot Q_{\rev,\mathrm{gen}}
=-I_{\rm rxn}T
\left.\frac{\partial U_{\eq}}{\partial T}\right|_{\rm state}.
}
\label{THM-003-HEAT}
\end{equation}
이 부호는 충전/방전 명칭이 아니라 식~\eqref{THM-003-CONVENTION}의
같은 forward reaction에서 따라온다.

\paragraph{\physid{THM-004} full-cell 식.}
forward cell reaction의 진행을 양으로 세어
$I_\cell=zF\dot n_\cell>0$로 고정하면
\begin{equation}
U_\cell=U_p-U_n,\qquad
\dot Q_{\rev,\mathrm{gen}}
=-I_\cell T\frac{\partial U_\cell}{\partial T}.
\label{THM-004-EQUATION}
\end{equation}
half-cell entropy coefficient를 어느 전극, 전해질 또는 계면의
local heat로 배정할지는 control volume과 reference reaction을
명시하기 전까지 열린 문제다.

\paragraph{\physid{THM-005} entropy basis의 택일.}
가역열은 다음 둘 중 하나의 basis로 계산한다.
\begin{enumerate}[label=(\Alph*)]
\item 측정되거나 유도된 fixed-state/fixed-charge OCV coefficient,
\item standard reaction entropy, configurational/partial-molar
entropy와 background-storage entropy를 모두 정의한 분해.
\end{enumerate}
두 basis를 독립 열원처럼 더하지 않는다. B를 택하면 그 합이
식~\eqref{THM-003-ENTROPY} 또는
식~\eqref{THM-002-EQUATION}와 같음을 유도해야 한다.
경험적 비대칭 형상계수는 이 entropy basis에 들어가지 않는다.

\subsection{비가역 entropy production}\label{THM-CH2-IRREVERSIBLE}

\paragraph{\physid{THM-006} local flux--affinity 법칙.}
반응 $j$의 molar flux와 chemical affinity가 공액이면
\begin{equation}
T\dot S_{\irr,j}
=\dot n_j\mathcal A_j^\chem\ge0.
\label{THM-006-LOCAL}
\end{equation}
두 상태 network에서는 $J_j^\pm>0$인 영역에서
\begin{equation}
\boxed{
T\dot S_{\irr,j}
=\frac{Q_j^{\rm SI}}{z_jF}RT
(J_j^+-J_j^-)
\ln\frac{J_j^+}{J_j^-}\ge0.
}
\label{THM-006-NETWORK}
\end{equation}
$Q_j^{\rm SI}$는 transition 전체의 C 단위 전하다. 따라서
$Q_j^{\rm SI}/(z_jF)$는 mol, $RT$는 J mol$^{-1}$,
$J_j^\pm$는 s$^{-1}$이어서 식의 단위는 W다.

비음성은 두 양수 $x,y$에 대해
$(x-y)\ln(x/y)\ge0$인 사실에서 직접 따른다. 식에 들어가는
$z_j$는 반응 화학양론이고 관측 폭의 함수가 아니다.

\paragraph{\physid{THM-007} terminal lumped approximation.}
하나의 signed terminal path에서 convention이
$I(U_{\rm oc}-V)\ge0$를 보장한다면
\begin{equation}
\dot Q_{\irr,\rm terminal}=I(U_{\rm oc}-V)
\label{THM-007-EQUATION}
\end{equation}
를 단자 손실의 lumped 근사로 쓸 수 있다. 이 식은 terminal
voltage gap에 포함된 소산만 세며, 휴지 중 내부 relaxation,
숨은 상태에너지 저장과 여러 local pathway를 포함하지 않는다.
식~\eqref{THM-006-NETWORK}의 일반 대체식이 아니다.

\paragraph{\physid{THM-008} transport 소산.}
물질 flux $\bm N_\ell$과 electrochemical-potential gradient가
공액이면 transport contribution은
\[
T\dot S_{\irr,\mathrm{tr}}
=\sum_\ell \bm N_\ell\boldsymbol{\cdot}
(-\nabla\widetilde\mu_\ell)\ge0
\]
로 둔다. 선형응답의 $I^2R$은 이 식의 제한된 축약이다. 관측
전압축의 이동계수를 곧바로 heat-generation resistance로
동일시하지 않는다.

\subsection{에너지 balance와 중복 계상}\label{THM-CH2-ENERGY}

\paragraph{\physid{BAL-005} control-volume 열수지.}
\begin{equation}
\boxed{
C_{\rm th}\dot T
=\dot Q_{\rev,\rm gen}
+T\dot S_{\irr,\rm rxn}
+T\dot S_{\irr,\rm tr}
-\dot Q_{\rm loss}
-\frac{\dd E_{\rm hidden}}{\dd t}.
}
\label{BAL-005-EQUATION}
\end{equation}
$E_{\rm hidden}$는 모형에 명시된 내부에너지 저장항이다.
이를 0으로 놓을 때만 모든 free-energy decrease를 즉시 열원으로
동일시할 수 있다.

\paragraph{\physid{THM-009} relaxation heat 중복 금지.}
같은 free-energy decrease를
$T\dot S_\irr$와 별도 ``relaxation heat''로 두 번 더하지 않는다.
가역 state-energy 변화, 저장되는 부분과 소산되는 부분을
식~\eqref{BAL-005-EQUATION} 하나에서 분해한다. 특정 tail
잔차가 $e^{-2t/\tau}$로 보인다는 사실만으로 power의 진폭이나
$1/\tau$ 계수를 정하지 않는다.

\begin{identifiability}
$\partial_TU$, heat loss, transport heat, reaction heat와 hidden
energy storage는 표면온도 한 신호에서 서로 보상할 수 있다.
준평형 entropy 측정, calorimetry, 전류차단과 공간적으로 정의된
control volume이 없으면 local heat allocation을 고유하게
식별하지 않는다.
\end{identifiability}

\begin{openitem}
\physid{THM-OPEN-01}
half-cell entropy heat의 전극/전해질/계면 공간 배분은 full
control-volume model과 calorimetry가 필요하다.
\end{openitem}
